% --- Discussion / Conclusion ---
% Science research articles end with a short discussion / outlook paragraph
% rather than a separate Conclusion section. Keep this tight (~250-400 words).

% Paragraph A — Central scientific takeaway.
%   - The active-variant region for a de novo enzyme is broader and farther
%     from the starting design than SSM or vanilla MPNN can reach.
%   - Function-aligned RL sampling + combinatorial library construction +
%     experimental closed loop together close the gap.
\todo{conclusion p1 — central takeaway tying the three applications together}

% Paragraph B — What each application individually established.
%   - Heme: fine-tuned MPNN dominates SSM and vanilla MPNN under matched
%     experimental effort — the foundation.
%   - PETase: experimental data → surrogate reward → improved next library;
%     a closed loop, not a single shot.
%   - Protease: explicit diversity pressure (batch-consensus divergence
%     reward) unlocks rare mutations and improves activity over an
%     otherwise-matched library.
\todo{conclusion p2 — what each app proves about the platform}

% Paragraph C — Honest limitations.
%   - PETase: probe and dimer activity demonstrated; bulk PET activity
%     remains an open challenge attributed to substrate accessibility on
%     solid plastic.
%   - Closed-loop iteration depends on having an assay-ready experimental
%     pipeline; cold-start is still hard.
%   - Diversity reward improves access to rare mutations but does not
%     substitute for a well-specified function reward.
\todo{conclusion p3 — limitations, including bulk PET gap}

% Paragraph D — Outlook.
%   - Generalization to additional chemistries and substrate classes
%     (including bulk polymers).
%   - Tighter integration of design-time rewards with experimental
%     surrogate signals across more rounds.
\todo{conclusion p4 — outlook}
